\documentclass[../../main.tex]{subfiles}% 注意这里的文件路径不能用 ./main.tex ,否则用latexmk编译子文件会报错
\graphicspath{{\subfix{./image/}}{\subfix{../../image/}}} % 指定图片引用路径,后续可以直接使用图片文件名
% 如果图片存放在上级目录,则使用 ../../image/ 作为备用路径

% 例如:
% \begin{figure}[H]
% \centering
% \includegraphics[width=0.3\textwidth]{图.png}
% \caption{}
% \label{figure:图}
% \end{figure}
% 注意:上述\label{}一定要放在\caption{}之后,否则引用图片序号会只会显示??.

\begin{document}

\section{场论}

\begin{definition}
设$D\subseteq\mathbb{R}^n$为非空开集，称映射
\begin{align*}
f:D\to\mathbb{R},\,\,\bm{x}=(x_1,\cdots,x_n)\mapsto f(\bm{x})
\end{align*}
为$D$上的\textbf{标量场}。
称映射
\begin{align*}
F:D\to\mathbb{R}^m,\,\,\bm{x}=(x_1,\cdots,x_n)\mapsto \bm{F}(\bm{x})
\end{align*}
为$D$上的\textbf{向量场}。通常取$m=n$，即
\begin{align*}
F:D\to\mathbb{R}^n,\,\,\bm{x}=(x_1,\cdots,x_n)\mapsto \bm{F}(\bm{x}).
\end{align*}
\end{definition}

\begin{definition}
设$f:D\subseteq\mathbb{R}^n\to\mathbb{R}$在点$\bm{x}_0$处所有一阶偏导数存在，称向量
\[
\nabla f(\bm{x}_0)=\left(\frac{\partial f}{\partial x_1}(\bm{x}_0),\cdots,\frac{\partial f}{\partial x_n}(\bm{x}_0)\right)
\]
为$f$在$\bm{x}_0$处的\textbf{梯度},称$\nabla$为$\mathbf{Hamilton}$\textbf{(哈密顿)算子}.
\end{definition}

\begin{proposition}\label{proposition:连通开集上梯度为0就是常数}
设 $\Omega\subset \mathbb{R}^n$ 是非空连通开集，$f:\Omega\to\mathbb{R}$ 在 $\Omega$ 上可微，且
\begin{align*}
\nabla f(x)=0,\qquad \forall x\in\Omega,
\end{align*}
则 $f$ 在 $\Omega$ 上为常数.
\end{proposition}
\begin{remark}
该命题的逆命题显然成立：常函数的梯度处处为零.
\end{remark}
\begin{proof}
任取 $x_0\in\Omega$，定义集合
\begin{align*}
A\triangleq\{\,x\in\Omega\mid f(x)=f(x_0)\,\}.
\end{align*}
显然$x_0\in A$，故$A\neq\varnothing$.

先证 $A$ 是开集：任取 $y\in A$，因 $\Omega$ 开，故存在 $r>0$ 使得 $B(y,r)\subset\Omega$. 对任意 $z\in B(y,r)$，由于球是凸的，线段 $[y,z]\subset B(y,r)\subset\Omega$. 对 $t\in[0,1]$，令 $\gamma(t)=y+t(z-y)$，则由链式法则得
\begin{align*}
\frac{\mathrm{d}}{\mathrm{d}t} f(\gamma(t))=\nabla f(\gamma(t))\cdot(z-y)=0.
\end{align*}
故 $f(\gamma(t))$ 恒为常数，因此$f(\gamma(0))=f(\gamma(1))=f(z)=f(y)=f(x_0)$，所以 $z\in A$. 因此 $B(y,r)\subset A$，即$A$ 为开集.

再证 $A$ 是闭集：由 $f$ 的连续性知$A=f^{-1}(\{f(x_0)\})$ 是 $\Omega$ 的闭集.

若$A \subsetneq  \Omega$,则$\Omega =A\sqcup (\Omega\setminus A)$,由$A$ 非空且既开又闭知$\Omega\setminus A$也是开集,这与$\Omega$ 连通矛盾!故 $A=\Omega$. 因此对任意 $x\in\Omega$，$f(x)=f(x_0)$，即 $f$ 为常数.
\end{proof}

\begin{proposition}\label{proposition:}
设 $A$ 是 $n$ 阶实矩阵，$x=(x_1,\cdots,x_n)^T\in\mathbb{R}^n$，定义二次型
\begin{align*}
f(x)\triangleq x^T A x=\sum_{i=1}^n\sum_{j=1}^n a_{ij}x_i x_j.
\end{align*}
则 $f$ 在 $x$ 处的梯度为
\begin{align}
\nabla f(x)=(A+A^T)x. \label{grad:general}
\end{align}
特别地，若 $A$ 为对称矩阵，则
\begin{align}
\nabla f(x)=2Ax. \label{grad:symmetric}
\end{align}
\end{proposition}
\begin{remark}
二次型 $x^T A x$ 的值仅由 $A$ 的对称部分 $\frac{1}{2}(A+A^T)$ 决定，因此梯度的表达式自然地反映了这一事实.
\end{remark}
\begin{proof}
将 $f(x)$ 展开：
\begin{align*}
f(x)=\sum_{i=1}^n\sum_{j=1}^n a_{ij}x_i x_j.
\end{align*}
固定 $k$，对 $x_k$ 求偏导：
\begin{align*}
\frac{\partial f}{\partial x_k}
=\sum_{i=1}^n\sum_{j=1}^n a_{ij}\frac{\partial}{\partial x_k}(x_i x_j) 
=\sum_{i=1}^n\sum_{j=1}^n a_{ij}(\delta_{ik}x_j+x_i\delta_{jk}) 
=\sum_{j=1}^n a_{kj}x_j+\sum_{i=1}^n a_{ik}x_i.
\end{align*}
上式第一项为 $(Ax)_k$，第二项为 $(A^T x)_k$，因此
\begin{align*}
\frac{\partial f}{\partial x_k}=((A+A^T)x)_k.
\end{align*}
所以 $\nabla f(x)=(A+A^T)x$，即 \eqref{grad:general}式成立.

当 $A$ 对称时，$A^T=A$，代入 \eqref{grad:general} 得 $\nabla f(x)=2Ax$，即 \eqref{grad:symmetric}式成立.
\end{proof}

\begin{definition}
设$\bm{F}=(F_1,\cdots,F_n)$是$\mathbb{R}^n$上的向量场，若$F_i$均可微，则定义$\bm{F}$的\textbf{散度}为
\[
\nabla\cdot\bm{F}=\sum_{i=1}^{n}\frac{\partial F_i}{\partial x_i}.
\]
\end{definition}

\begin{definition}
设$\bm{F}=(F_1,F_2,F_3)$是$\mathbb{R}^3$上的向量场，若各分量均可微，则定义$\bm{F}$的\textbf{旋度}为
\[
\nabla\times\bm{F}=
\begin{vmatrix}
\bm{i} & \bm{j} & \bm{k} \\
\frac{\partial}{\partial x} & \frac{\partial}{\partial y} & \frac{\partial}{\partial z} \\
F_1 & F_2 & F_3
\end{vmatrix}
=\left(\frac{\partial F_3}{\partial y}-\frac{\partial F_2}{\partial z},\frac{\partial F_1}{\partial z}-\frac{\partial F_3}{\partial x},\frac{\partial F_2}{\partial x}-\frac{\partial F_1}{\partial y}\right).
\]
\end{definition}

\begin{definition}
对于二次可微的标量场$f:\mathbb{R}^n\to\mathbb{R}$，定义\textbf{Laplace(拉普拉斯)算子}为
\[
\Delta f=\nabla\cdot(\nabla f)=\sum_{i=1}^{n}\frac{\partial^{2}f}{\partial x_i^{2}}.
\]
\end{definition}

\begin{definition}
设$D\subseteq\mathbb{R}^n$为开集，$f:D\to\mathbb{R}$，$\bm{x}_0\in D$，$\bm{e}\in\mathbb{R}^n$为单位向量。若极限
\begin{align*}
\lim_{t\to 0}\frac{f(\bm{x}_0+t\bm{e})-f(\bm{x}_0)}{t}
\end{align*}
存在，则称该极限为$f$在点$\bm{x}_0$处沿方向$\bm{e}$的\textbf{方向导数}，记作$\frac{\partial f}{\partial\bm{e}}(\bm{x}_0)$或$D_{\bm{e}}f(\bm{x}_0)$.
\end{definition}

\begin{theorem}[方向导数的基本性质]\label{theorem:方向导数的基本性质}
设$f,g:D\subseteq\mathbb{R}^n\to\mathbb{R}$在点$\bm{x}_0\in D$处可微，$\alpha,\beta\in\mathbb{R}$，$\bm{e}\in \mathbb{R}^n$为任意单位向量，$\nabla f(\bm{x}_0)$为梯度。则
\begin{enumerate}[(1)]
\item\label{theorem:方向导数的基本性质-1} \textbf{梯度表示:}方向导数存在且
\begin{align*}
\frac{\partial f}{\partial\bm{e}}(\bm{x}_0)=\nabla f(\bm{x}_0)\cdot\bm{e}.
\end{align*}

\item \textbf{线性性质:}
\begin{align*}
\frac{\partial(\alpha f+\beta g)}{\partial\bm{e}}(\bm{x}_0)=\alpha\frac{\partial f}{\partial\bm{e}}(\bm{x}_0)+\beta\frac{\partial g}{\partial\bm{e}}(\bm{x}_0).
\end{align*}

\item 若$\nabla f(\bm{x}_0)\neq\bm{0}$，则当$\bm{e}$与$\nabla f(\bm{x}_0)$同向时方向导数取得最大值$|\nabla f(\bm{x}_0)|$，反向时取得最小值$-|\nabla f(\bm{x}_0)|$.进而对所有单位方向$\bm{e}$有
\begin{align*}
-|\nabla f(\bm{x}_0)|\leqslant\frac{\partial f}{\partial\bm{e}}(\bm{x}_0)\leqslant|\nabla f(\bm{x}_0)|.
\end{align*}

\item 设$\varphi(t)=f(\bm{x}_0+t\bm{e})$，则$\varphi'(0)=\frac{\partial f}{\partial\bm{e}}(\bm{x}_0)$.并且当$f$可微时,有$\varphi'(t)=\nabla f(\bm{x}_0+t\bm{e})\cdot\bm{e}$.
\end{enumerate}
\end{theorem}
\begin{remark}
方向导数存在不能推出函数在该点连续，即使所有方向导数存在且为零，函数仍可能不连续。例如
\[
f(x,y)=\begin{cases}\dfrac{x^{2}y}{x^{4}+y^{2}},&(x,y)\neq(0,0),\\0,&(x,y)=(0,0),\end{cases}
\]
在原点沿任何方向的方向导数均为$0$，但$f$在原点不连续.
\end{remark}

\begin{theorem}\label{theorem:散度旋度Laplace算子的性质}
设$f,g$是标量场，$\bm{F},\bm{G}$是向量场，$c\in \mathbb{R}$，则
\begin{enumerate}[(1)]
\item\label{theorem:散度旋度Laplace算子的性质-1} $\nabla(cf)=c\nabla f$，$\quad \nabla\cdot(c\bm{F})=c\nabla\cdot\bm{F}$，$\quad \nabla\times(c\bm{F})=c\nabla\times\bm{F}$.

\item\label{theorem:散度旋度Laplace算子的性质-2} $\nabla(f\pm g)=\nabla f\pm\nabla g$，$\quad \nabla\cdot(\bm{F}\pm\bm{G})=\nabla\cdot\bm{F}\pm\nabla\cdot\bm{G}$，$\quad \nabla\times(\bm{F}\pm\bm{G})=\nabla\times\bm{F}\pm\nabla\times\bm{G}$.

\item\label{theorem:散度旋度Laplace算子的性质-3} $\nabla(fg)=f\nabla g+g\nabla f$.

\item\label{theorem:散度旋度Laplace算子的性质-4} $\nabla\cdot(f\bm{F})=f\nabla\cdot\bm{F}+\nabla f\cdot\bm{F}$.

\item\label{theorem:散度旋度Laplace算子的性质-5} $\nabla\times(f\bm{F})=f\nabla\times\bm{F}+\nabla f\times\bm{F}$.

\item\label{theorem:散度旋度Laplace算子的性质-6} $\nabla\cdot(\bm{F}\times\bm{G})=\bm{G}\cdot(\nabla\times\bm{F})-\bm{F}\cdot(\nabla\times\bm{G})$.

\item\label{theorem:散度旋度Laplace算子的性质-7} 旋度无源:$\nabla\times(\nabla f)=\bm{0}$.

\item\label{theorem:散度旋度Laplace算子的性质-8} 散度无旋:$\nabla\cdot(\nabla\times\bm{F})=0$.

\item\label{theorem:散度旋度Laplace算子的性质-9} $\Delta(cu)=c\Delta u$
\item $\Delta(u\pm v)=\Delta u\pm\Delta v$.

\item\label{theorem:散度旋度Laplace算子的性质-10} $\Delta(uv)=u\Delta v+v\Delta u+2\nabla u\cdot\nabla v$,特别地$\Delta u^2=2u\Delta u+\left| \nabla u \right|^2.$
\end{enumerate}
\end{theorem}

\begin{theorem}[Gauss-Green公式(散度定理)]\label{theorem:Gauss-Green公式(散度定理)}
设$\Omega$ 是 $\mathbb{R}^n$ 中的一个有界开集，且 $\partial\Omega\in C^1$。若向量场 $\bm{F}=(F_1,F_2,\cdots,F_n):\Omega\to\mathbb{R}^n$且$\bm{F}\in C^1(\Omega)\cap C(\overline{\Omega})$，则
\begin{align*}
\int_\Omega \operatorname{div} \bm{F}\,\mathrm{d}x =\int_{\Omega}\nabla\cdot\bm{F}\,\mathrm{d}V= \int_{\partial\Omega} \bm{F}\cdot \bm{n}\,\mathrm{d}S(x).
\end{align*}
其中 $\bm{n}$ 为 $\partial\Omega$ 的单位外法向量。
\end{theorem}

\begin{theorem}[Green公式（平面散度定理与旋度定理）]\label{theorem:Green公式平面散度定理与旋度定理}
设$D\subset\mathbb{R}^2$是有界闭区域，边界$\partial D$是分段光滑的闭曲线，取正向（逆时针）。若$P,Q$在$\overline{D}$上连续，在$D$内连续可微，则
\begin{align*}
\iint_{D}\left(\frac{\partial Q}{\partial x}-\frac{\partial P}{\partial y}\right)\mathrm{d}x\mathrm{d}y=\oint_{\partial D}P\,\mathrm{d}x+Q\,\mathrm{d}y, 
\end{align*}
\begin{align}
\iint_{D}\left(\frac{\partial P}{\partial x}+\frac{\partial Q}{\partial y}\right)\mathrm{d}x\mathrm{d}y=\oint_{\partial D}(P\,\mathrm{d}y-Q\,\mathrm{d}x).\label{eq:green_div} 
\end{align}
\end{theorem}
\begin{remark}
上述定理中的\eqref{eq:green_div}是散度定理在二维的体现.
\end{remark}

\begin{theorem}[Green恒等式]\label{theorem:Green恒等式}
设$\Omega\subset\mathbb{R}^n$是有界闭区域，其边界$\partial\Omega$是逐片光滑的定向曲面，$\bm{n}$为$\partial\Omega$的单位外法向量。
\begin{enumerate}[(1)]
\item\label{theorem:Green第一恒等式} $\mathbf{Green}$\textbf{第一恒等式:}若$u,v\in C^2(\Omega)\cap C^1(\overline{\Omega})$，则有
\begin{align*}
\int_{\Omega}u\Delta v\,\mathrm{d}V=\int_{\partial\Omega}u\frac{\partial v}{\partial\bm{n}}\,\mathrm{d}S-\int_{\Omega}\nabla u\cdot\nabla v\,\mathrm{d}V.
\end{align*}
特别地,取$u\equiv 1$,则有
\begin{align*}
\int_{\Omega}\Delta v\,\mathrm{d}V=\int_{\partial\Omega}\frac{\partial v}{\partial\bm{n}}\,\mathrm{d}S.
\end{align*}

\item\label{theorem:Green第二恒等式} $\mathbf{Green}$\textbf{第二恒等式:}若$u,v\in C^2(\Omega)\cap C^1(\overline{\Omega})$，则有
\begin{align*}
\int_{\Omega}(u\Delta v-v\Delta u)\,\mathrm{d}V=\int_{\partial\Omega}\left(u\frac{\partial v}{\partial\bm{n}}-v\frac{\partial u}{\partial\bm{n}}\right)\mathrm{d}S. 
\end{align*}

\item\label{theorem:Green恒等式-3} 若$u,v\in C^1(\Omega)\cap C(\overline{\Omega})$，则有
\begin{align*}
\int_\Omega \frac{\partial u}{\partial x_i}v\,\mathrm{d}x = -\int_\Omega u \frac{\partial v}{\partial x_i}\,\mathrm{d}x + \int_{\partial\Omega} u v n_i\,\mathrm{d}S(x),
\end{align*}
其中 $n_i$ 是 $\partial\Omega$ 的单位外法向量 $\bm{n}$ 的第 $i$ 个分量.特别地，取$v\equiv 1$，则有
\begin{align*}
\int_{\Omega}\frac{\partial u}{\partial x_i}\,\mathrm{d}V=\int_{\partial\Omega}u\,n_i\,\mathrm{d}S. 
\end{align*}
\end{enumerate}
\end{theorem}
\begin{proof}
\begin{enumerate}[(1)]
\item 取 $\bm{F} = u \nabla v$，其中 $u, v \in C^2(\Omega) \cap C^1(\overline{\Omega})$。则由\rrefthe{theorem:散度旋度Laplace算子的性质}{theorem:散度旋度Laplace算子的性质-4}和Laplace算子的定义知
\[
\nabla \cdot \bm{F} = \nabla \cdot (u \nabla v) = \nabla u \cdot \nabla v + u \nabla \cdot (\nabla v) = \nabla u \cdot \nabla v + u \Delta v.
\]
又由$u, \nabla v \in C^1(\overline{\Omega})$知$\bm{F} \in C^1(\overline{\Omega};\mathbb{R}^n)$，故由\hyperref[theorem:Gauss-Green公式(散度定理)]{散度定理}和\hyperref[theorem:方向导数的基本性质-1]{方向导数的基本性质}得
\[
\int_{\Omega}\nabla\cdot\bm{F}\,\mathrm{d}V=\int_{\partial\Omega}\bm{F}\cdot\bm{n}\,\mathrm{d}S \iff \int_{\Omega} (u \Delta v + \nabla u \cdot \nabla v) \, \mathrm{d}V = \int_{\partial\Omega} (u \nabla v) \cdot \bm{n} \, \mathrm{d}S = \int_{\partial\Omega} u \frac{\partial v}{\partial \bm{n}} \, \mathrm{d}S,
\]
其中 $\frac{\partial v}{\partial \bm{n}} = \nabla v \cdot \bm{n}$ 是 $v$ 沿外法向$\bm{n}$的方向导数。移项即得
\begin{align}\label{eq::HJIOUJWFIOJO3424288}
\int_{\Omega} u \Delta v \, \mathrm{d}V = \int_{\partial\Omega} u \frac{\partial v}{\partial \bm{n}} \, \mathrm{d}S - \int_{\Omega} \nabla u \cdot \nabla v \, \mathrm{d}V.
\end{align}

\item 将\eqref{eq::HJIOUJWFIOJO3424288}式中的 $u$ 与 $v$ 互换，得到
\[
\int_{\Omega} v \Delta u \, \mathrm{d}V = \int_{\partial\Omega} v \frac{\partial u}{\partial \bm{n}} \, \mathrm{d}S - \int_{\Omega} \nabla v \cdot \nabla u \, \mathrm{d}V.
\]
将\eqref{eq::HJIOUJWFIOJO3424288}式减去上式，再结合$\nabla u \cdot \nabla v = \nabla v \cdot \nabla u$，便有
\[
\int_{\Omega} (u \Delta v - v \Delta u) \, \mathrm{d}V = \int_{\partial\Omega} \left( u \frac{\partial v}{\partial \bm{n}} - v \frac{\partial u}{\partial \bm{n}} \right) \mathrm{d}S.
\]

\item 取向量场
\begin{align*}
\bm{F}=\left( 0,\cdots, 0, uv, 0, \cdots, 0 \right),
\end{align*}
其中非零分量在第$i$个位置, 即$\bm{F}=uv\bm{e}_i$. 由\hyperref[theorem:Gauss-Green公式(散度定理)]{散度定理}得
\begin{align*}
\int_{\Omega}{\nabla \cdot \mathbf{F}\,\mathrm{d}x}=\int_{\partial \Omega}{\mathbf{F}\cdot \mathbf{n}\,\mathrm{d}S}&\Longleftrightarrow \int_{\Omega}{\frac{\partial}{\partial x_i}\left( uv \right) \mathrm{d}x}=\int_{\Omega}{\left( \frac{\partial u}{\partial x_i}v+u\frac{\partial v}{\partial x_i} \right) \mathrm{d}x}=\int_{\partial \Omega}{uvn_i\,\mathrm{d}S}
\\
&\Longleftrightarrow \int_{\Omega}{\frac{\partial u}{\partial x_i}v\mathrm{d}x}=-\int_{\Omega}{u\frac{\partial v}{\partial x_i}\mathrm{d}x}+\int_{\partial \Omega}{uvn_i\,\mathrm{d}S}.
\end{align*}
\end{enumerate}
\end{proof}

\begin{theorem}[多元函数分部积分公式]\label{theorem:多元函数分部积分公式}
设$\Omega\subset\mathbb{R}^n$是有界闭区域，其边界$\partial\Omega$是逐片光滑的定向曲面，$\bm{n}$为$\partial\Omega$的单位外法向量。若$u\in C^1(\Omega)\cap C(\overline{\Omega})$，$\bm{F}\in C^1(\Omega)\cap C(\overline{\Omega})$为向量场，则
\begin{align*}
\int_{\Omega}u(\nabla\cdot\bm{F})\,\mathrm{d}V=\int_{\partial\Omega}u(\bm{F}\cdot\bm{n})\,\mathrm{d}S-\int_{\Omega}\nabla u\cdot\bm{F}\,\mathrm{d}V. 
\end{align*}
\end{theorem}

























\end{document}